% Replot Figure 8(b) with the original curves and measured radius samples.
% Build from paper/figures after running code/export_figure8b_data.py.
\documentclass[tikz,border=0pt]{standalone}
\usepackage{fontspec}
\usepackage{xeCJK}
\usepackage{pgfplots}
\setmainfont{Times New Roman}
\setCJKmainfont{SimSun}
\pgfplotsset{compat=1.18}
\definecolor{RadiusGreen}{HTML}{009E73}
\definecolor{MeanBlue}{HTML}{56B4E9}
\definecolor{ReferenceGray}{HTML}{666666}
\def\DryingEndHour{51.090529435454968}

\pgfplotsset{
  original panel/.style={
    scale only axis,
    width=228.247788bp,height=101.988679bp,
    xmin=0,xmax=72,
    xtick={0,20,40,60},xticklabels={0,20,40,60},
    axis lines=box,
    axis line style={black,line width=0.8bp},
    tick align=outside,
    xtick pos=bottom,ytick pos=left,
    tick style={black,line width=0.8bp},
    major tick length=3.5bp,
    tick label style={font=\fontsize{10.5}{12}\selectfont},
    label style={font=\fontsize{11.4}{13}\selectfont},
    grid=major,
    grid style={draw=black!15,line width=0.5bp},
    axis background/.style={fill=white},
    scaled ticks=false,
    clip=true,
    clip marker paths=true,
  },
}
\begin{document}
\begin{tikzpicture}[x=1bp,y=1bp]
  % The former 600 dpi export was 2391 x 2298 px. Keep that same canvas.
  \path[use as bounding box] (0,0) rectangle (286.92,275.76);
  \begin{pgfinterruptboundingbox}
    \begin{axis}[
      original panel,
      at={(51.450962bp,151.633821bp)},anchor=south west,
      ymin=1.1579,ymax=2.0401,
      ytick={1.25,1.50,1.75,2.00},
      yticklabels={1.25,1.50,1.75,2.00},
      xticklabels={},
      ylabel={半径(cm)},
      ylabel style={at={(axis description cs:-0.145,0.5)},anchor=south},
      legend style={
        at={(0.985,0.975)},anchor=north east,
        fill=white,fill opacity=1,text opacity=1,draw=none,
        font=\fontsize{8.4}{9.7}\selectfont,
        inner sep=1bp,row sep=0.5bp,
      },
      % Shorter samples keep the entire white legend to the right of the endpoint.
      legend image code/.code={
        \draw[#1,mark repeat=2,mark phase=2]
          plot coordinates {(0bp,0bp) (5bp,0bp) (10bp,0bp)};
      },
      legend cell align=left,
    ]
      \addplot[ReferenceGray,line width=1.1bp,densely dotted,forget plot]
        coordinates {(\DryingEndHour,1.1579) (\DryingEndHour,2.0401)};
      \addplot[only marks,mark=o,mark size=1.6bp,
        mark options={draw=black!80,fill=white,line width=0.8bp}]
        table[x=time,y=radius] {fig08b_measurements.dat};
      \addlegendentry{附件2实测}
      \addplot[RadiusGreen,line width=1.8bp]
        table[x=time,y=radius] {fig08b_curves.dat};
      \addlegendentry{插值 R(t)}
      \node[anchor=north east,align=center,inner sep=1bp,fill=white,
        font=\fontsize{8.5}{10}\selectfont]
        at (axis cs:49.5,1.75) {烘干结束\\51.09小时};
    \end{axis}
    \node[anchor=base west,inner sep=0pt,
      font=\fontsize{11.4}{13}\selectfont]
      at (51.450962,261.0) {(b) 半径收缩与平均含水率};
    \begin{axis}[
      original panel,
      at={(51.450962bp,37.4065bp)},anchor=south west,
      ymin=-0.0026577835112786,ymax=2.660457670556573,
      ytick={0,1,2},yticklabels={0,1,2},
      ylabel={平均含水率(kg/kg)},
      ylabel style={at={(axis description cs:-0.145,0.5)},anchor=south},
      xlabel={时间(h)},
      xlabel style={at={(axis description cs:0.5,-0.13)},anchor=north},
    ]
      \addplot[ReferenceGray,line width=1.1bp,densely dotted]
        coordinates {(\DryingEndHour,-0.0026577835112786)
          (\DryingEndHour,2.660457670556573)};
      \addplot[MeanBlue,line width=1.8bp]
        table[x=time,y=mean] {fig08b_curves.dat};
      \node[anchor=north east,align=center,inner sep=1bp,fill=white,
        font=\fontsize{8.5}{10}\selectfont]
        at (axis cs:49.5,1.30) {烘干结束\\51.09小时};
    \end{axis}
  \end{pgfinterruptboundingbox}
\end{tikzpicture}
\end{document}
